\section{Introduction --- draft v1}
\label{sec:intro}

During this compiler's development, its semantic analyzer detected a
program defect and told no one. The program was `int main() { int x
= 5; }` --- a non-void function that never returns. The analyzer's
internal log recorded the failure precisely ("bare return in
non-void function"); no diagnostic was ever constructed from that
knowledge; compilation "succeeded"; and the produced binary returned
whatever value happened to occupy the return register. The compiler
knew. The user didn't. We found the defect not by reading the code
but by running twenty adversarial programs against the finished
compiler --- and we report it as the paper's motivating observation
because it is the general phenomenon in miniature: **the evidence
behind compiler diagnostics is routinely discarded before, or
instead of, reaching the user.** What a conventional pipeline
retains at the point of failure --- a message string and a source
location --- is a small projection of what it knew moments earlier:
the construct being analyzed, the rule being applied, the path taken
to get there, and what became of the user's code in every stage
downstream.

This is not the familiar complaint that error messages are unhelpful.
Message \emph{quality} has received sustained attention in production
compilers --- Rust's structured suggestions \cite{rust-rfc1644}, GHC's migration to
errors-as-values \cite{ghc-structured-errors} [$\rightarrow$\S5] --- and a substantial research literature
examines message effectiveness, with notably contested results:
controlled studies of \emph{enhanced} messages report effects ranging
from ineffectual to inconclusive \cite{denny2014enhancing,pettit2017enhanced}. We take that contestation seriously: this
paper makes no claim that its diagnostics help users more, and the
user study such a claim requires is designed but not run (\S13). Our
question is upstream of wording, and architectural: **what must a
compiler retain, and where must it retain it, for explanation to be
possible at all --- and what does retaining it cost?**

Two scope commitments frame everything that follows. First, the
vehicle is a deliberately small compiler: a C-like kernel language
(integers, arithmetic, functions without calls) compiled through a
complete pipeline --- lexing, recursive-descent parsing, semantic
analysis, three-address IR, three optimization passes run to a fixed
point, and 32-bit x86 code generation to a linked executable. The
contribution is the architecture and its measured costs, not a
production technique; scaling questions are discussed, not claimed
away (\S10). Second, "explanation" here means *preserved compiler
evidence* --- recorded execution paths, provenance, root causes --- not
post-hoc rationalization, and not machine-learning explainability.

The architecture rests on one obligation: no stage may discard the
evidence a future explanation would need. Concretely (Fig. F1 shows
one diagnostic carrying all of it): every stage returns its product
\emph{and} its diagnostics, so partial results survive errors; a
diagnostic is a structured value --- kind, severity, span, root cause,
explanation, fixes, examples, trace --- never a string; all prose
lives in one knowledge base indexed by root cause, so the taxonomy
is the explanation index; each stage owns a trace recorder whose
RAII scopes capture the \emph{actual} execution path, snapshotted at the
moment a diagnostic is created; every IR instruction carries its
source span and accumulates notes from the optimizer passes that
transform it, surfacing as source-line comments in the emitted
assembly; and the properties that make diagnostics trustworthy ---
one mistake yields one diagnostic; traces contain runtime facts ---
are regression tests, not aspirations.

Making that architecture honest required measuring it, and the
measurements pushed back. Our first overhead campaign measured the
harness, not the compiler, and was discarded (\S9.2 explains how, so
readers do not repeat it). The second refuted our own hypothesis:
naive trace recording --- eagerly building detail strings at every
scope entry --- cost 2.4$\times$ compile time at 20k lines, dominated by two
scopes per token in the lexer. The repair (lazy detail formatting
over caller-owned data, exploiting the fact that snapshots occur
only while frames are open) brought recording to 1--2\% with
confidence intervals spanning zero, with byte-identical output.
Provenance, by contrast, is \emph{not} free: a steady 12--17\%, mechanical
in origin --- the span and note fields grow every IR instruction from
116 to 160 bytes, and every copy and traversal pays. We also report
what provenance cannot promise: per-instruction history is not
transitive under dead-code elimination (a folded instruction's
history dies with it once propagation makes it dead), and our
cascade-freedom invariants, while holding on every tested shape,
leave a measured residual --- 3\% of single-fault mutants still
produce three or more diagnostics.

This paper contributes:

\begin{itemize}
\item \textbf{C1 --- an architectural pattern for explanation-oriented compilation:} the five-role diagnostic core (root-cause taxonomy, factory, isolated knowledge base, multi-target rendering) over a pipeline-wide output-plus-diagnostics contract --- presented as the distilled, designed-in form of a structure production compilers (GHC, Rust) have been converging toward by retrofit \cite{rust-rfc1644,ghc-structured-errors} [$\rightarrow$\S5].
\item \textbf{C2 --- execution-recorded diagnostic traces:} RAII scope recording with snapshot-at-creation and a curated fallback that is provably unreachable in the pipeline; with the lazy-formatting design that makes always-on recording statistically free, and the measurements for both designs.
\item \textbf{C3 --- end-to-end provenance under the same contract:} source spans and transformation notes surviving to emitted assembly --- conceding LLVM's optimization remarks \cite{llvm-remarks} as production-scale prior art for the mechanism, and contributing its unification with the diagnostic architecture, its measured price, and the DCE non-transitivity finding.
\item \textbf{C4 --- diagnostic quality as executable invariants:} property- level tests (cascade-freedom, trace accuracy) distinct from golden-output snapshots, with mutation testing quantifying the frontier they do not yet cover.
\end{itemize}

\S4 fixes vocabulary and the language; \S5 positions against five
literature clusters; \S6 presents the architecture as fourteen
design decisions, each argued against a rejected alternative; \S7
describes the visualization consumer; \S8 the implementation; \S9 the
three-iteration measurement campaign and quality studies; Sections~10--13
discuss scaling, validity, limitations, and future work, including
the designed user study; \S14 concludes.



